\documentclass[11pt,a4paper]{article}

% ── Packages ──────────────────────────────────────────────────────────────
\usepackage[margin=1in]{geometry}
\usepackage{amsmath,amssymb}
\usepackage{booktabs}
\usepackage{longtable}
\usepackage{multirow}
\usepackage{array}
\usepackage{xcolor}
\usepackage{hyperref}
\usepackage{url}
\usepackage{graphicx}
\usepackage{enumitem}
\usepackage{natbib}
\usepackage{microtype}
\usepackage{caption}
\usepackage[T1]{fontenc}
\usepackage[utf8]{inputenc}

\hypersetup{
  colorlinks=true,
  linkcolor=blue!60!black,
  citecolor=teal!70!black,
  urlcolor=blue!60!black
}

% ── Custom commands ────────────────────────────────────────────────────────
\newcommand{\metricname}[1]{\textbf{#1}}
\newcommand{\abbr}[1]{\texttt{#1}}
\newcommand{\arxiv}[1]{\href{https://arxiv.org/abs/#1}{arXiv:#1}}

% ── Title ──────────────────────────────────────────────────────────────────
\title{%
  \textbf{A Comprehensive Survey of Evaluation Metrics for}\\
  \textbf{LLM-based Multi-Agent Systems}\\[6pt]
  \large Taxonomy, Classification, and Practical Guidance
}

\author{%
  [Author Name(s)]\\
  [Institution(s)]\\
  \texttt{[email@institution.edu]}
}

\date{\today}

% ==========================================================================
\begin{document}
\maketitle

% ── Abstract ───────────────────────────────────────────────────────────────
\begin{abstract}
The rapid deployment of Large Language Model (LLM)-based multi-agent systems (MAS)
in research and industry has exposed a critical gap: evaluation methodologies designed
for single-agent systems fail to capture the emergent, collective behaviours that arise
when multiple agents interact. This survey consolidates \textbf{95 evaluation metrics}
drawn from 61 primary research papers, organising them into a nine-category taxonomy
spanning Task Performance, Communication Quality, Coordination and Collaboration,
Safety and Trust, Planning and Reasoning, Medical/Domain-Specific, Framework and
Governance, Operational Infrastructure, and Tool Use. We additionally introduce a
principled classification of metrics into \emph{trajectory-level} (54 metrics requiring
observation of the full execution sequence) and \emph{outcome-level} (41 metrics
requiring only the final output). For each metric we report the formal definition,
value range, directionality, source benchmark or paper, and a practitioner importance
rating on a five-point scale. We conclude by identifying the most significant open
challenges in MAS evaluation: the absence of standardised privacy leakage metrics,
the under-measurement of hallucination propagation across agent boundaries,
coordination deadlock quantification, and run-to-run reproducibility. This survey
serves as a reference guide for researchers designing MAS evaluation protocols
and practitioners selecting metrics for production deployments.
\end{abstract}

\tableofcontents
\newpage

% ==========================================================================
\section{Introduction}
\label{sec:intro}

Multi-agent systems built on large language models are increasingly deployed across
high-stakes domains including software engineering \citep{swe_bench2024}, scientific
research, healthcare \citep{medaide2024}, and enterprise automation. Unlike classical
MAS, LLM-based agents can coordinate via natural language, use external tools, plan
over long horizons, and adapt their behaviour dynamically. These capabilities introduce
evaluation challenges that go well beyond measuring whether a single agent produces the
correct final output.

Three properties make LLM-MAS evaluation qualitatively different from single-agent
evaluation:

\begin{enumerate}[leftmargin=*, label=\textbf{\arabic*.}]
  \item \textbf{Emergence.} System-level behaviours—including coordinated bias
    amplification \citep{emergent_bias2024} and hallucination propagation across agent
    boundaries—cannot be predicted from the individual agent scores.
  \item \textbf{Trajectory dependence.} The quality of a MAS run often depends not on
    whether a task was ultimately completed, but on \emph{how} it was completed: the
    communication efficiency, the division of labour, the recovery strategy after
    failures, and the safety of intermediate steps.
  \item \textbf{Framework sensitivity.} Empirical studies show that the choice of
    orchestration framework (e.g., AutoGen \citep{autogen2023}, LangGraph, Claude SDK)
    alone causes performance ranges comparable to differences between frontier
    models \citep{maseval2025}.
\end{enumerate}

Despite these challenges, no comprehensive reference consolidates the full landscape
of proposed metrics. Individual papers introduce metrics for narrow settings, making it
difficult for practitioners to select appropriate evaluation protocols. This survey fills
that gap.

\paragraph{Contributions.}
\begin{itemize}
  \item A unified taxonomy of 95 metrics across 9 categories, with formal definitions,
    value ranges, and directionality for each.
  \item A novel trajectory-level versus outcome-level classification that guides metric
    selection based on what aspects of system behaviour are being studied.
  \item A practitioner importance rating (1--5) grounded in deployment frequency and
    regulatory relevance.
  \item An analysis of open evaluation gaps: unmeasured metrics, reproducibility
    concerns, and emergent MAS-specific phenomena.
\end{itemize}

% ==========================================================================
\section{Background and Related Work}
\label{sec:background}

\subsection{LLM Agent Evaluation Surveys}

Recent surveys have catalogued evaluation approaches for individual LLM agents.
\citet{llm_eval_survey2025} organise evaluation along two dimensions: what to evaluate
(behaviour, capabilities, reliability, safety) and how to evaluate (interaction modes,
benchmarks, metric computation). \citet{agent_eval_survey2025} survey evaluation
frameworks for tool-using agents with a focus on robustness and safety. Our work
differs in scope: we focus exclusively on \emph{multi-agent} systems and cover both
research and production-grade operational metrics.

\subsection{Key Benchmarks}

Several benchmarks have driven MAS evaluation:

\begin{itemize}
  \item \textbf{GAIA} \citep{gaia2023}: 466 general-assistant tasks graded by Goal
    Success Rate (GSR). The de-facto standard for comparing MAS framework performance.
  \item \textbf{$\tau$-bench} \citep{tbench2024}: 120 tool-agent-user interaction tasks
    graded by \abbr{pass@k} and \abbr{pass\^{}k}, distinguishing lenient from strict
    reliability.
  \item \textbf{MMLU} \citep{mmlu2020}: 14,042 tasks measuring task accuracy across
    57 academic subjects; widely used as a capability baseline.
  \item \textbf{MultiAgentBench (MARBLE)} \citep{marble2025}: Evaluates collaboration
    and competition across diverse agent topologies with milestone-based KPIs.
  \item \textbf{AgentBoard} \citep{agentboard2024}: Multi-turn evaluation with
    fine-grained progress tracking across 9 agent task categories.
  \item \textbf{SWE-bench} \citep{swe_bench2024}: 500 human-verified GitHub issues;
    the primary benchmark for code-generation MAS.
\end{itemize}

\subsection{Agent Frameworks Under Evaluation}

Empirical comparisons in \citet{maseval2025} cover AutoGen \citep{autogen2023},
LangGraph, Claude SDK, smolagents, CAMEL \citep{camel2023}, and LlamaIndex.
The GAIA GSR scores for these frameworks range from 44.6\% to 53.2\%, a spread nearly
equal to the 49.7\%--55.4\% range observed across frontier models
(GPT-4o, Claude 3.5 Sonnet, Gemini 1.5 Pro).

% ==========================================================================
\section{Metric Taxonomy}
\label{sec:taxonomy}

We organise 95 metrics into the taxonomy shown in Table~\ref{tab:taxonomy_summary}.
The nine categories reflect distinct aspects of MAS behaviour, from end-to-end task
outcomes to per-token operational costs.

\begin{table}[h!]
  \centering
  \caption{Summary of the 95-metric taxonomy.}
  \label{tab:taxonomy_summary}
  \begin{tabular}{llcc}
    \toprule
    \textbf{Category} & \textbf{Focus} & \textbf{Count} & \textbf{Traj.} \\
    \midrule
    Task Performance         & End-to-end task outcomes            & 14 & 2  \\
    Communication Quality    & Inter-agent message efficiency       & 10 & 10 \\
    Coordination             & Role allocation, game-theoretic      & 20 & 15 \\
    Safety \& Trust          & Attacks, alignment, reliability      & 23 & 13 \\
    Planning \& Reasoning    & Intermediate steps, memory           &  7 &  7 \\
    Medical / Domain         & Domain-specific output quality       &  5 &  0 \\
    Framework \& Governance  & Systemic and fairness properties     &  7 &  3 \\
    Operational              & Latency, cost, reproducibility       &  6 &  4 \\
    Tool Use                 & Per-call tool accuracy               &  3 &  3 \\
    \midrule
    \textbf{Total}           &                                     & \textbf{95} & \textbf{54} \\
    \bottomrule
  \end{tabular}
\end{table}

\noindent\textbf{Trajectory-level vs.\ outcome-level.} A metric is
\emph{trajectory-level} if it requires observing the full execution sequence of agent
actions—i.e., it cannot be computed from only the final output. \emph{Outcome-level}
metrics require only the final result. Of the 95 metrics, 54 are trajectory-level and
41 are outcome-level.

% ==========================================================================
\section{Task Performance Metrics}
\label{sec:task}

Task performance metrics quantify whether agents complete assigned goals, with varying
degrees of credit for partial progress.

\subsection{Goal Success Rate (GSR)}

\begin{equation}
  \text{GSR} = \frac{\text{Completed Tasks}}{\text{Total Tasks}} \times 100\%
\end{equation}

GSR is the universal MAS leaderboard metric \citep{maseval2025,gaia2023}. Framework
choice alone causes an 8.6 percentage-point GSR range on GAIA. \textbf{Importance: 5
(Critical)}.

\subsection{pass@k and pass\^{}k}

Introduced in $\tau$-bench \citep{tbench2024}, these two metrics measure fundamentally
different reliability properties:

\begin{align}
  \text{pass@}k &= 1 - \frac{\binom{n-c}{k}}{\binom{n}{k}} \quad
    \text{(probability that at least one of } k \text{ trials succeeds)} \\[4pt]
  \text{pass}^k &= \left(\frac{c}{n}\right)^k \quad
    \text{(probability that all } k \text{ trials succeed)}
\end{align}

\noindent where $n$ is the total number of trials and $c$ is the number of successes.
\abbr{pass\^{}k} drops exponentially with $k$ and is the strict reliability metric
favoured for production evaluation.

\subsection{Partial Goal Success Rate (pGSR)}

\begin{equation}
  \text{pGSR} = \frac{\sum_{i} w_i \cdot s_i}{\sum_{i} w_i}
\end{equation}

pGSR awards partial credit for sub-goal completion, making it more informative than
binary GSR for long-horizon tasks \citep{maseval2025}.

\subsection{Additional Task Metrics}

Table~\ref{tab:task_metrics} summarises the remaining task performance metrics.

\begin{table}[h!]
  \centering
  \caption{Additional task performance metrics.}
  \label{tab:task_metrics}
  \small
  \begin{tabular}{lllc}
    \toprule
    \textbf{Abbr.} & \textbf{Name} & \textbf{Source} & \textbf{Range} \\
    \midrule
    Sub.R    & Subgoal Success Rate      & \citet{comm_verify2025}  & 0--100\%  \\
    Sup-GSR  & Supervisor GSR            & \citet{aws_bedrock2024}  & 0--100\%  \\
    Sys-GSR  & System-side GSR           & \citet{aws_bedrock2024}  & 0--100\%  \\
    BCS      & Behavior Competency Score & \citet{crew_wildfire2025}& 0--1      \\
    WR       & Win Rate                  & \citet{magic2023}        & 0--100\%  \\
    KPI      & Key Performance Indicator & \citet{marble2025}       & task-def. \\
    PassRate & Pass Rate                 & \citet{mobile_bench2024} & 0--100\%  \\
    SuccR    & Success Rate              & \citet{safe_beal2025}    & 0--100\%  \\
    SafeR@S  & Safety Rate at Success    & \citet{safe_beal2025}    & 0--100\%  \\
    \bottomrule
  \end{tabular}
\end{table}

% ==========================================================================
\section{Communication Quality Metrics}
\label{sec:comm}

Communication metrics evaluate the efficiency and quality of information exchange
between agents. All are trajectory-level: they require the full message trace.

\subsection{Information Divergence Score (IDS)}

\begin{equation}
  \text{IDS} = \frac{\sum_{(i,j)} w_{ij}\cdot(1 - \text{SS}_{\text{total}}[i,j])}
                    {\sum_{(i,j)} w_{ij}}
\end{equation}

where $w_{ij} = \max(S_{ij}, S_{ji}) + \max(T_{ij}, T_{ji})$ and
$\text{SS}_{\text{total}} = 0.5 \cdot \text{SS}_{\text{syn}} + 0.5 \cdot
\text{SS}_{\text{sem}}$ (syntactic and semantic similarity). Introduced in
GEMMAS \citep{gemmas2025}, IDS revealed a 12.8\% divergence gap between two systems
with near-identical GSM8K accuracy, demonstrating that outcome metrics can mask
communication quality differences.

\subsection{Unnecessary Path Ratio (UPR)}

\begin{equation}
  \text{UPR} = 1 - \frac{|P_{\text{necessary}}|}{|P_{\text{all}}|}
\end{equation}

UPR quantifies redundant reasoning paths in the agent communication DAG. On GSM8K,
\citet{gemmas2025} found an 80\% UPR gap between two systems with otherwise similar
accuracy, illustrating how outcome metrics obscure efficiency differences.

\subsection{Team Execution Score (TES)}

\begin{equation}
  \text{TES}(\bar{h}_k) = \max_j \left\{
    \frac{(1+\beta^2) \cdot D_{\max}^j(\bar{h}_k, \bar{g}_k^j)}{m_k + \beta^2 n_k}
  \right\}
\end{equation}

where $\bar{h}_k$ is the agent history for subtask $k$, $\bar{g}_k^j$ is the $j$-th
ground-truth trajectory, $m_k$ counts false-positive actions, and $n_k$ counts
false-negative actions. \citet{collab_overcooked2025} found that LLMs degrade sharply
at complexity level 4+, particularly in initiating capability (IC) versus responding
capability (RC).

\subsection{Additional Communication Metrics}

\begin{table}[h!]
  \centering
  \caption{Additional communication quality metrics.}
  \label{tab:comm_metrics}
  \small
  \begin{tabular}{lllc}
    \toprule
    \textbf{Abbr.} & \textbf{Name} & \textbf{Source} & \textbf{Range} \\
    \midrule
    IC   & Initiating Capability          & \citet{collab_overcooked2025} & 0--100\% \\
    RC   & Responding Capability          & \citet{collab_overcooked2025} & 0--100\% \\
    ITES & Incremental Team Exec.\ Score  & \citet{collab_overcooked2025} & 0--1     \\
    PC   & Progress Completeness          & \citet{collab_overcooked2025} & 0--1     \\
    AF   & Alignment Factor               & \citet{c2c2025}               & 0--1     \\
    CORE & Conv.\ Robustness Eval.\ Score & CORE (2025)                   & 0--1     \\
    NCCG & Norm.\ Cumulative Common Ground& \citet{collab_deliberate2025} & 0--1     \\
    \bottomrule
  \end{tabular}
\end{table}

\noindent The Alignment Factor (AF) from the C2C framework \citep{c2c2025} is
noteworthy: systems optimising AF reduced task completion time by approximately 40\%
with acceptable cost increases across teams of 5--17 agents.

% ==========================================================================
\section{Coordination and Collaboration Metrics}
\label{sec:coord}

Coordination metrics evaluate how effectively agents allocate roles, synchronise
actions, and resolve conflicts.

\subsection{Average Normalised Utility (ANU)}

ANU \citep[COLLAB;][]{marble2025} measures the aggregate utility of coordination
actions normalised by the theoretical maximum. It captures whether agents coordinate
to increase the total value produced, not merely to complete tasks.

\subsection{Milestone Completion Rate (MCR)}

\begin{equation}
  \text{MCR} = \frac{\text{Milestones Completed}}{\text{Total Milestones}} \times 100\%
\end{equation}

Introduced in MultiAgentBench \citep{marble2025}, MCR awards partial credit that binary
GSR does not. A cognitive planning strategy increased MCR by 3 percentage points, while
GSR remained unchanged—demonstrating that MCR reveals information that outcome metrics
miss.

\subsection{Division of Labor Index (DoLI)}

\begin{equation}
  \text{DoLI} = 1 - \frac{H(\text{task distribution})}{\log N_{\text{agents}}}
\end{equation}

where $H$ is Shannon entropy over the task assignment distribution. DoLI = 1 indicates
perfect specialisation; DoLI = 0 indicates all agents perform identical work.
\citet{llm_eval_survey2025} identify DoLI as ``almost never measured'' despite being
fundamental to MAS efficiency.

\subsection{Social Welfare Score (SWS)}

\begin{equation}
  \text{SWS} = \sum_{i=1}^{N} R_i
\end{equation}

where $R_i$ is the cumulative reward of agent $i$. SWS is a game-theoretic metric that
captures whether the MAS achieves collectively optimal outcomes versus individually
rational but collectively suboptimal Nash equilibria \citep{game_theoretic_mas2025}.

\subsection{MAgIC Social Behaviour Metrics}

The MAgIC benchmark \citep{magic2023} decomposes agent social behaviour across five
dimensions measured during game-theoretic interactions:

\begin{equation}
  \{S_J,\, S_R,\, S_D,\, S_{\text{self}},\, S_{\text{collab}}\}
\end{equation}

covering Judgment, Reasoning, Deception, Self-Awareness, and Cooperation respectively.

\subsection{Additional Coordination Metrics}

\begin{table}[h!]
  \centering
  \caption{Coordination and collaboration metrics.}
  \label{tab:coord_metrics}
  \small
  \begin{tabular}{lllc}
    \toprule
    \textbf{Abbr.} & \textbf{Name} & \textbf{Source} & \textbf{Range} \\
    \midrule
    CS               & Coordination Score          & \citet{marble2025}          & 0--1      \\
    Elo              & Elo Rating                  & \citet{magic2023}           & open      \\
    FASR             & False Agent Switch Rate     & \citet{aws_bedrock2024}     & 0--100\%  \\
    Conformity Rate  & Conformity Rate             & \citet{benchform2025}       & 0--100\%  \\
    Independence Rate& Independence Rate           & \citet{benchform2025}       & 0--100\%  \\
    CoM Rate         & Change-of-Mind Rate         & \citet{collab_deliberate2025}& 0--100\% \\
    CSS              & Component Synergy Score     & \citet{trism2025}           & 0--1      \\
    TUE              & Tool Utilisation Efficacy   & \citet{trism2025}           & 0--1      \\
    CoS              & Collaboration Score         & MindAgent                   & 0--1      \\
    CompS            & Competition Score           & \citet{marble2025}          & 0--1      \\
    NSR              & Negotiation Success Rate    & \citet{marble2025}          & 0--100\%  \\
    \bottomrule
  \end{tabular}
\end{table}

% ==========================================================================
\section{Safety and Trust Metrics}
\label{sec:safety}

Safety metrics form the largest category (23 metrics), reflecting the diverse threat
surface of deployed MAS.

\subsection{Attack Success Rate (ASR)}

\begin{equation}
  \text{ASR} = \frac{\text{Successful Attacks}}{\text{Total Attack Attempts}} \times 100\%
\end{equation}

ASR is the universal adversarial safety metric \citep{tamas2025}. It is the complement
of the defender's success rate and is used across Agent SafetyBench and AutoDefense.
\textbf{Importance: 5 (Critical)}.

\subsection{Over-Exposure Rate (OER) and Alignment Dispersion (AD)}

From the Trust-Vulnerability Paradox study \citep{trust_vuln2025}, OER measures the
fraction of runs in which agents disclose more information than authorised. AD measures
the statistical dispersion of safety scores across runs:

\begin{equation}
  \text{AD} = \text{Var}\!\left(\text{safety\_score}\right) \text{ across N runs}
\end{equation}

A high AD score with acceptable mean safety is particularly dangerous in production:
the system is \emph{unpredictably} safe.

\subsection{Effective Robustness Score (ERS)}

\begin{equation}
  \text{ERS} = \frac{\text{Tasks completed safely under adversarial conditions}}
                    {\text{Total tasks}} \times 100\%
\end{equation}

\citet{tamas2025} show that ERS and ASR can diverge significantly: a system may
successfully defend against attacks (low ASR) while still failing to complete tasks
safely (low ERS).

\subsection{Over-Refuse Rate (ORR)}

\begin{equation}
  \text{ORR}(D_{\text{borderline}}, \pi) = \frac{1}{N}\sum_{i=1}^{N}
  \mathbf{1}[\text{benign prompt}_i \text{ refused}] \times 100\%
\end{equation}

ORR from WaltzRL \citep{waltzrl2025} measures the safety-helpfulness trade-off. Overly
cautious agents reduce task utility; ORR and ASR together define the Pareto frontier.

\subsection{PII Leakage Rate (PLR)}

\begin{equation}
  \text{PLR} = \frac{\text{Leaked PII instances}}
                    {\text{Total sensitive instances}} \times 100\%
\end{equation}

AgentLeak \citep{agentleak2025} is the first full-stack MAS privacy benchmark (1,000
scenarios; healthcare, finance, legal, corporate). A critical MAS-specific finding:
multi-agent configurations \emph{reduce} per-channel output leakage (27.2\% vs. 43.2\%
single-agent) but \emph{increase} total system exposure to 68.9\% because internal
agent-to-agent channels are unmonitored.

\subsection{Hallucination Propagation Rate (HPR)}

\begin{equation}
  \text{HPR} = \frac{\text{Outputs containing upstream hallucination}}
                    {\text{Total outputs}} \times 100\%
\end{equation}

HPR is a uniquely MAS metric \citep{cleanlab2025}: a single hallucinating agent in a
sequential topology contaminates all downstream consumers. No single-agent metric can
detect this phenomenon. Trust scoring reduces hallucination-propagated errors by 55.8\%
in ReAct-type agents.

\subsection{Ethical Cooperation Score (ECS)}

\begin{equation}
  \text{ECS} = C \times A \times I \times F
\end{equation}

where $C$, $A$, $I$, $F$ are Cooperation, Autonomy, Integrity, and Fairness scores.
The multiplicative form from CMAG \citep{cmag2025} ensures that any single dimension
approaching zero collapses the overall score, penalising manipulative behaviour.

\subsection{Additional Safety Metrics}

\begin{table}[h!]
  \centering
  \caption{Additional safety and trust metrics.}
  \label{tab:safety_metrics}
  \small
  \begin{tabular}{lllc}
    \toprule
    \textbf{Abbr.} & \textbf{Name} & \textbf{Source} & \textbf{Range} \\
    \midrule
    PDR      & Process Danger Rate          & \citet{psysafe2024}      & 0--100\% \\
    ASV      & Avg. Safety Violations       & \citet{safer_robotics2025}& 0--1    \\
    Robustness & Robustness Score           & \citet{maseval2025}      & 0--1     \\
    FPR      & False Positive Rate          & \citet{guardians2025}    & 0--100\% \\
    JDR      & Joint Danger Rate            & \citet{psysafe2024}      & 0--100\% \\
    PNA      & Performance under No Attack  & \citet{tamas2025}        & 0--100\% \\
    SafeR    & Safety Rate                  & \citet{safe_beal2025}    & 0--100\% \\
    APDR     & Avg. Performance Drop Rate   & \citet{pptcr2024}        & 0--100\% \\
    FTR      & Feedback Trigger Rate        & \citet{waltzrl2025}      & 0--100\% \\
    Autonomy & Autonomy Score               & \citet{cmag2025}         & 0--1     \\
    DR       & Deadlock Rate                & \citet{llm_eval_survey2025}& 0--100\%\\
    TAR      & Task Abandonment Rate        & \citet{llm_eval_survey2025}& 0--100\%\\
    RR       & Recovery Rate                & \citet{llm_eval_survey2025}& 0--100\%\\
    HR       & Hallucination Rate           & \citet{cleanlab2025}     & 0--100\% \\
    ECE      & Expected Calibration Error   & \citet{cleanlab2025}     & 0--1     \\
    \bottomrule
  \end{tabular}
\end{table}

% ==========================================================================
\section{Planning and Reasoning Metrics}
\label{sec:planning}

All seven planning metrics are trajectory-level, requiring observation of intermediate
reasoning and action steps.

\subsection{CheckPoint Score and Progress Rate}

CheckPoint \citep{mobile_bench2024} measures intermediate goal completion during
execution. Progress Rate \citep{agentboard2024} provides a fine-grained view of
sub-goal achievement at each step:

\begin{equation}
  \text{Progress Rate} = \frac{\text{Sub-goals achieved at step } t}
                               {\text{Total sub-goals}} \times 100\%
\end{equation}

\subsection{Step Ratio (StepR)}

\begin{equation}
  \text{StepR} = \frac{\text{Executed steps}}{\text{Optimal solution steps}}
\end{equation}

StepR from \citet{comm_verify2025} measures planning efficiency. Values greater than 1
indicate unnecessary steps; values less than 1 indicate incomplete execution.

\subsection{Memory and Long-horizon Metrics}

\metricname{Memory Recall Accuracy (MRA)} measures how accurately agents retrieve
relevant prior context across multi-turn interactions \citep{llm_eval_survey2025}.
Typical LLMs achieve MRA $\approx 92\%$ over the last 5 turns but only $\approx 41\%$
over 50+ turns without external memory augmentation.

\metricname{Long-Horizon Coherence (LHC)} measures consistency of stated facts,
intent, and behaviour across 10+ agent steps:

\begin{equation}
  \text{LHC} = 1 - \text{contradiction\_rate}(\text{steps}_{1..N})
\end{equation}

\subsection{Reasoning Chain Metrics}

Average Reasoning Length (ARL) and CoT Match Rate from Structured Agent Distillation
\citep{structured_distill2025} measure the length and step-by-step fidelity of
chain-of-thought reasoning chains against ground-truth trajectories.

% ==========================================================================
\section{Medical and Domain-Specific Metrics}
\label{sec:medical}

Five metrics address domain-specific output quality in medical and clinical MAS:

\begin{itemize}
  \item \metricname{AUROC} and \metricname{AUPRC}: Standard classification quality
    metrics applied to clinical decision tasks in MedAgentBoard.
  \item \metricname{ROUGE-L}: Longest common subsequence-based recall for clinical
    text generation.
  \item \metricname{BLEU}: N-gram precision for medical report generation
    \citep{medaide2024}.
  \item \metricname{SARI}: System output against references and input; specialised for
    text simplification tasks in patient-facing MAS.
\end{itemize}

% ==========================================================================
\section{Framework and Governance Metrics}
\label{sec:framework}

\subsection{Bias Amplification Factor (BAF)}

\begin{equation}
  \text{BAF} = \frac{\text{bias\_score}(\text{MAS output})}
                    {\text{bias\_score}(\text{single agent})}
\end{equation}

\citet{emergent_bias2024} demonstrate that a MAS can exhibit group-preference bias even
when no individual agent exhibits prior bias—a purely emergent phenomenon. BAF $> 1$
indicates MAS-amplified bias; BAF $< 1$ indicates MAS-corrected bias. EU AI Act and
US SR 11-7 require bias audits for high-stakes deployments.

\subsection{Agent Scaling Efficiency (ASE)}

\begin{equation}
  \text{ASE}_{\text{norm}} = \frac{\text{performance}_N}{N \times \text{performance}_1}
\end{equation}

where $N$ is agent count. Most MAS show $\text{ASE}_{\text{norm}} \approx 0.3$--$0.6$
beyond 4 agents, reflecting quadratic growth in coordination overhead (the Metcalfe
penalty). ASE drives capacity planning decisions.

\subsection{Trustworthiness Index (T) and Composite Metric ($M_{\text{comp}}$)}

From TRiSM \citep{trism2025}, $T$ aggregates trust signals across agent interactions,
while $M_{\text{comp}}$ provides a five-axis composite covering capability, robustness,
safety, human-centredness, and economic sustainability.

% ==========================================================================
\section{Operational and Infrastructure Metrics}
\label{sec:ops}

Operational metrics bridge the gap between research evaluation and production
deployment. All were absent from MAS research until 2025.

\subsection{Latency: TTFT and TPOT}

\begin{align}
  \text{TTFT} &= t_{\text{first token}} - t_{\text{request received}} \quad\text{(ms)} \\
  \text{TPOT} &= \frac{t_{\text{last token}} - t_{\text{first token}}}
                       {\text{output tokens} - 1} \quad\text{(ms/token)}
\end{align}

MLPerf 2025 production SLAs: TTFT $\leq 500$\,ms and TPOT $\leq 30$\,ms for
interactive agents \citep{nvidia_nim2025}. In a multi-agent pipeline, TTFT and TPOT
apply at each agent boundary; errors compound across hops.

\subsection{Cost Per Task (CPT)}

\begin{equation}
  \text{CPT} = \frac{\text{input\_tokens} \times p_{\text{in}} +
                     \text{output\_tokens} \times p_{\text{out}}}
                    {\text{Successful Tasks}} \quad\text{(USD)}
\end{equation}

CPT is the primary ROI metric for enterprise MAS deployment. GPT-4o MAS pipelines
can cost 5--50$\times$ more per task than single-agent baselines, and framework choice
alone causes 2--4$\times$ CPT variation \citep{galileo_benchmarks2025}.

\subsection{Goodput}

\begin{equation}
  \text{Goodput} = \frac{\text{Requests meeting SLA}}{\text{Total requests}} \times 100\%
\end{equation}

The 2025 consensus metric replacing raw throughput \citep{databricks_inference2025}:
a system at 1,000\,req/s with 60\% goodput is worse than 500\,req/s at 95\% goodput.

\subsection{Run-to-Run Variance (RunVar)}

\begin{equation}
  \text{RunVar} = \sigma^2(\text{GSR across } N \text{ identical runs})
\end{equation}

MAESTRO \citep{maestro2025} identified MAS architecture (not model temperature) as the
dominant driver of run-to-run variance. A system with GSR $= 0.70 \pm 0.30$ is
substantially less reliable than one with GSR $= 0.60 \pm 0.05$.

% ==========================================================================
\section{Tool Use Metrics}
\label{sec:tool}

Tool use metrics decompose end-to-end tool call quality into three orthogonal signals:

\begin{align}
  \text{TSA} &= \frac{\text{Correctly selected tools}}{\text{Total tool calls}} \times 100\%
              \quad\text{(Tool Selection Accuracy)} \\[4pt]
  \text{TPA} &= \frac{\text{Correct parameters}}{\text{Expected parameters}} \times 100\%
              \quad\text{(Tool Parameter Accuracy)} \\[4pt]
  \text{TCR} &= \frac{\text{Successful executions}}{\text{Total tool invocations}} \times 100\%
              \quad\text{(Tool Call Success Rate)}
\end{align}

\citet{llm_eval_survey2025} report that GPT-4o achieves TSA $\approx 95\%$ on simple
APIs but only $\approx 78\%$ on complex multi-tool workflows where options overlap
semantically. Crucially, high TSA does not imply high TPA: GPT-4o achieves TSA $\approx
95\%$ but TPA $\approx 82\%$, indicating that parameter specification is a distinct
failure mode from tool selection.

% ==========================================================================
\section{Trajectory-Level vs.\ Outcome-Level Classification}
\label{sec:trajectory}

We propose a principled classification that distinguishes how much of the execution
trace is required to compute a metric.

\begin{description}
  \item[Outcome-level (41 metrics)] Computable from a single input--output pair.
    Includes all final task success measures (GSR, Accuracy, ASR, pass@k) and
    aggregate quality scores (AUROC, BLEU, ROUGE-L). These are simpler to compute
    but miss all within-run dynamics.

  \item[Trajectory-level (54 metrics)] Require the full execution sequence of agent
    states, actions, and communications. Include all communication metrics (IDS, UPR,
    TES, IC, RC, ITES, PC, AF, CORE, NCCG), all planning metrics (Progress Rate,
    StepR, CheckPoint, ARL, CoT Match Rate, MRA, LHC), per-call tool metrics (TSA,
    TPA, TCR), per-step operational metrics (TTFT, TPOT, TBU), and event-accumulation
    safety metrics (OER, ERS, PDR, FPR, HPR, DR, PLR).
\end{description}

Table~\ref{tab:trajectory_examples} illustrates this distinction with representative
examples.

\begin{table}[h!]
  \centering
  \caption{Representative trajectory-level and outcome-level metric pairs.}
  \label{tab:trajectory_examples}
  \begin{tabular}{lll}
    \toprule
    \textbf{Aspect} & \textbf{Outcome-level} & \textbf{Trajectory-level} \\
    \midrule
    Task completion   & GSR                & Progress Rate, MCR     \\
    Communication     & --                 & IDS, UPR, TES          \\
    Safety            & ASR                & OER, HPR, DR, PLR      \\
    Planning          & --                 & StepR, ARL, LHC        \\
    Tool use          & --                 & TSA, TPA, TCR          \\
    Latency           & --                 & TTFT, TPOT             \\
    Bias              & BAF                & --                     \\
    \bottomrule
  \end{tabular}
\end{table}

% ==========================================================================
\section{Practitioner Importance Ratings}
\label{sec:importance}

We rate each metric's real-world importance on a five-point scale:

\begin{itemize}
  \item \textbf{5 -- Critical}: Industry-standard leaderboard metric; required for
    any publication or deployment. (GSR, Accuracy, ASR)
  \item \textbf{4 -- High}: Widely adopted in production-grade evaluations;
    expected in enterprise deployment reports. (26 metrics)
  \item \textbf{3 -- Medium}: Research standard; important for academic rigour.
    (47 metrics)
  \item \textbf{2 -- Emerging}: Novel (2024--25); important future direction but
    not yet standardised. (16 metrics)
  \item \textbf{1 -- Niche}: Highly specialised; relevant only for specific settings.
\end{itemize}

Table~\ref{tab:full_metrics} (Appendix~\ref{app:full_table}) provides the complete
95-metric reference with importance ratings, formulae, and citation keys.

% ==========================================================================
\section{Open Challenges and Research Gaps}
\label{sec:gaps}

Our survey reveals several significant evaluation gaps.

\subsection{Unmeasured MAS-Specific Phenomena}

\paragraph{Hallucination propagation.}
HPR has no existing benchmark implementation. The 2025 Cleanlab study \citep{cleanlab2025}
provides the first empirical evidence of cross-agent contamination, but no standardised
dataset or tool exists for measuring HPR.

\paragraph{Deadlock detection.}
DR is identified by \citet{llm_eval_survey2025} as ``almost never measured'' despite
being a production failure mode. Detecting deadlocks requires tracing agent wait
dependencies across a run—a capability absent from most evaluation frameworks.

\paragraph{Privacy through internal channels.}
\citet{agentleak2025} demonstrate that MAS internal channels are the primary PII
exposure vector, yet no standard evaluation protocol addresses this.

\subsection{Reproducibility}

\citet{maestro2025} show that MAS architecture is the dominant driver of run-to-run
variance. Academic papers rarely report RunVar or confidence intervals over multiple
runs, making results difficult to compare or reproduce.

\subsection{Framework vs.\ Model Attribution}

\citet{maseval2025} show that the GSR range across frameworks (44.6\%--53.2\%) nearly
equals the range across frontier models (49.7\%--55.4\%). Standard benchmark reporting
conflates these two sources of variance; disentangling them requires factorial designs.

\subsection{Scalability Evaluation}

ASE measurements in the literature show most MAS achieve $\text{ASE}_{\text{norm}}
\approx 0.3$--$0.6$ beyond 4 agents, yet optimal agent count is almost never reported
as part of benchmark results.

% ==========================================================================
\section{Recommended Practitioner Metric Stacks}
\label{sec:stacks}

Based on importance ratings and coverage analysis, we recommend the following
evaluation stacks:

\paragraph{Minimum viable stack (4 metrics).}
GSR + ASR + CPT + RunVar. Covers task success, safety, cost, and reliability.

\paragraph{Research standard stack (8 metrics).}
GSR + pass\^{}k + pGSR + ASR + ERS + OER + IDS + Progress Rate. Adds reliability,
partial credit, dual safety dimensions, and communication quality.

\paragraph{Production-grade stack (12 metrics).}
All of the above, plus: ORR + HR + TSA + TCR + TTFT + CPT. Adds
safety-helpfulness balance, hallucination, tool use quality, and operational SLAs.

\paragraph{Comprehensive evaluation (18+ metrics).}
Add MCR, DoLI, HPR, PLR, RunVar, BAF, ECE for full coverage of coordination,
privacy, calibration, and fairness.

% ==========================================================================
\section{Conclusion}
\label{sec:conclusion}

This survey consolidates 95 evaluation metrics for LLM-based multi-agent systems,
drawn from 61 primary research papers published between 2023 and 2025. Our nine-category
taxonomy and trajectory/outcome classification provide a structured framework for metric
selection. The three Critical-importance metrics (GSR, Accuracy, ASR) provide a
baseline; the 26 High-importance metrics form a production-grade evaluation stack. Key
open challenges—hallucination propagation, privacy leakage through internal channels,
deadlock detection, and reproducibility—represent the highest-priority research
directions for the field.

As LLM-MAS move from research settings to regulated deployment contexts (EU AI Act,
HIPAA, SR 11-7), the demand for standardised, auditable evaluation protocols will
intensify. We hope this survey accelerates convergence on such protocols.

% ==========================================================================
\bibliographystyle{plainnat}
\bibliography{references}

% ==========================================================================
\appendix
\section{Complete Metric Reference Table}
\label{app:full_table}

\begin{longtable}{p{1.8cm}p{3.5cm}p{1.2cm}p{1.8cm}p{1.0cm}c}
  \caption{All 95 evaluation metrics with category, source, range, directionality, and importance.}
  \label{tab:full_metrics} \\
  \toprule
  \textbf{Abbr.} & \textbf{Name} & \textbf{Cat.} & \textbf{Source} & \textbf{Range} & \textbf{Imp.} \\
  \midrule
  \endfirsthead
  \multicolumn{6}{c}{\textit{Table~\ref{tab:full_metrics} continued}} \\
  \toprule
  \textbf{Abbr.} & \textbf{Name} & \textbf{Cat.} & \textbf{Source} & \textbf{Range} & \textbf{Imp.} \\
  \midrule
  \endhead
  \bottomrule
  \endlastfoot
  % ── Task Performance ──
  GSR            & Goal Success Rate              & Task    & MASEval           & 0--100\%   & 5 \\
  Accuracy       & Task Accuracy                  & Task    & MMLU              & 0--100\%   & 5 \\
  pass@k         & pass@k (lenient)               & Task    & $\tau$-bench      & 0--100\%   & 4 \\
  pass\^{}k      & pass\^{}k (strict)             & Task    & $\tau$-bench      & 0--100\%   & 4 \\
  pGSR           & Partial GSR                    & Task    & MASEval           & 0--100\%   & 4 \\
  Sub.R          & Subgoal Success Rate           & Task    & CommVerify        & 0--100\%   & 4 \\
  Sup-GSR        & Supervisor GSR                 & Task    & AWS Bedrock       & 0--100\%   & 4 \\
  Sys-GSR        & System-side GSR                & Task    & AWS Bedrock       & 0--100\%   & 4 \\
  BCS            & Behavior Competency Score      & Task    & CREW-Wildfire     & 0--1       & 3 \\
  WR             & Win Rate                       & Task    & MAgIC             & 0--100\%   & 3 \\
  KPI            & Key Performance Indicator      & Task    & MARBLE            & task-def.  & 3 \\
  PassRate       & Pass Rate                      & Task    & Mobile-Bench      & 0--100\%   & 3 \\
  SuccR          & Success Rate                   & Task    & Safe-BeAl         & 0--100\%   & 3 \\
  SafeR@S        & Safety Rate at Success         & Task    & Safe-BeAl         & 0--100\%   & 3 \\
  \midrule
  % ── Communication ──
  IDS            & Information Divergence Score   & Comm.   & GEMMAS            & 0--1       & 3 \\
  UPR            & Unnecessary Path Ratio         & Comm.   & GEMMAS            & 0--1       & 3 \\
  TES            & Team Execution Score           & Comm.   & Collab-OC         & 0--1       & 3 \\
  IC             & Initiating Capability          & Comm.   & Collab-OC         & 0--100\%   & 3 \\
  RC             & Responding Capability          & Comm.   & Collab-OC         & 0--100\%   & 3 \\
  AF             & Alignment Factor               & Comm.   & C2C               & 0--1       & 2 \\
  CORE           & Conv.\ Robustness Eval.        & Comm.   & CORE              & 0--1       & 2 \\
  ITES           & Incr.\ Team Exec.\ Score       & Comm.   & Collab-OC         & 0--1       & 2 \\
  PC             & Progress Completeness          & Comm.   & Collab-OC         & 0--1       & 2 \\
  NCCG           & Norm.\ Cumul.\ Common Ground   & Comm.   & Collab+Del.       & 0--1       & 2 \\
  \midrule
  % ── Coordination ──
  ANU            & Avg.\ Normalised Utility       & Coord.  & COLLAB            & 0--1       & 3 \\
  CS             & Coordination Score             & Coord.  & MARBLE            & 0--1       & 3 \\
  Elo            & Elo Rating                     & Coord.  & MAgIC             & open       & 4 \\
  FASR           & False Agent Switch Rate        & Coord.  & AWS Bedrock       & 0--100\%   & 3 \\
  Conformity     & Conformity Rate                & Coord.  & BenchForm         & 0--100\%   & 3 \\
  Independence   & Independence Rate              & Coord.  & BenchForm         & 0--100\%   & 3 \\
  CoM Rate       & Change-of-Mind Rate            & Coord.  & Collab+Del.       & 0--100\%   & 3 \\
  CSS            & Component Synergy Score        & Coord.  & TRiSM             & 0--1       & 2 \\
  TUE            & Tool Utilisation Efficacy      & Coord.  & TRiSM             & 0--1       & 3 \\
  S\_J           & Judgment (MAgIC)               & Coord.  & MAgIC             & 0--1       & 2 \\
  S\_R           & Reasoning (MAgIC)              & Coord.  & MAgIC             & 0--1       & 3 \\
  S\_D           & Deception (MAgIC)              & Coord.  & MAgIC             & 0--1       & 2 \\
  S\_self        & Self-Awareness (MAgIC)         & Coord.  & MAgIC             & 0--1       & 2 \\
  S\_collab      & Cooperation (MAgIC)            & Coord.  & MAgIC             & 0--1       & 3 \\
  CoS            & Collaboration Score            & Coord.  & MindAgent         & 0--1       & 2 \\
  MCR            & Milestone Completion Rate      & Coord.  & MARBLE            & 0--100\%   & 4 \\
  CompS          & Competition Score              & Coord.  & MARBLE            & 0--1       & 3 \\
  DoLI           & Division of Labor Index        & Coord.  & Survey            & 0--1       & 3 \\
  NSR            & Negotiation Success Rate       & Coord.  & MARBLE            & 0--100\%   & 2 \\
  SWS            & Social Welfare Score           & Coord.  & Game-Theory MAS   & $>0$       & 3 \\
  \midrule
  % ── Safety ──
  ASR            & Attack Success Rate            & Safety  & TAMAS             & 0--100\%   & 5 \\
  OER            & Over-Exposure Rate             & Safety  & Trust-Vuln.       & 0--100\%   & 4 \\
  AD             & Alignment Dispersion           & Safety  & Trust-Vuln.       & $\geq 0$   & 4 \\
  ERS            & Effective Robustness Score     & Safety  & TAMAS             & 0--100\%   & 4 \\
  ORR            & Over-Refuse Rate               & Safety  & WaltzRL           & 0--100\%   & 4 \\
  PDR            & Process Danger Rate            & Safety  & PsySafe           & 0--100\%   & 4 \\
  ASV            & Avg.\ Safety Violations        & Safety  & SAFER             & 0--1       & 4 \\
  Robustness     & Robustness Score               & Safety  & MASEval           & 0--1       & 4 \\
  FPR            & False Positive Rate            & Safety  & Guardians         & 0--100\%   & 4 \\
  ECS            & Ethical Cooperation Score      & Safety  & CMAG              & 0--1       & 3 \\
  JDR            & Joint Danger Rate              & Safety  & PsySafe           & 0--100\%   & 3 \\
  PNA            & Perf.\ under No Attack         & Safety  & TAMAS             & 0--100\%   & 3 \\
  SafeR          & Safety Rate                    & Safety  & Safe-BeAl         & 0--100\%   & 3 \\
  APDR           & Avg.\ Performance Drop Rate    & Safety  & PPTC-R            & 0--100\%   & 3 \\
  FTR            & Feedback Trigger Rate          & Safety  & WaltzRL           & 0--100\%   & 3 \\
  Autonomy       & Autonomy Score                 & Safety  & CMAG              & 0--1       & 2 \\
  DR             & Deadlock Rate                  & Safety  & Survey            & 0--100\%   & 3 \\
  TAR            & Task Abandonment Rate          & Safety  & Survey            & 0--100\%   & 3 \\
  RR             & Recovery Rate                  & Safety  & Survey            & 0--100\%   & 3 \\
  HR             & Hallucination Rate             & Safety  & Cleanlab          & 0--100\%   & 4 \\
  HPR            & Hallucin.\ Propagation Rate    & Safety  & Cleanlab          & 0--100\%   & 3 \\
  ECE            & Expected Calibration Error     & Safety  & Cleanlab          & 0--1       & 3 \\
  PLR            & PII Leakage Rate               & Safety  & AgentLeak         & 0--100\%   & 4 \\
  \midrule
  % ── Planning ──
  CheckPoint     & CheckPoint Score               & Plan.   & Mobile-Bench      & 0--1       & 4 \\
  Progress Rate  & Fine-Grained Progress Rate     & Plan.   & AgentBoard        & 0--100\%   & 4 \\
  StepR          & Step Ratio                     & Plan.   & CommVerify        & $>0$       & 3 \\
  ARL            & Avg.\ Reasoning Length         & Plan.   & Struct.\ Distill. & steps      & 3 \\
  CoT Match Rate & CoT Match Rate                 & Plan.   & Struct.\ Distill. & 0--100\%   & 3 \\
  MRA            & Memory Recall Accuracy         & Plan.   & Survey            & 0--100\%   & 3 \\
  LHC            & Long-Horizon Coherence         & Plan.   & Survey            & 0--1       & 2 \\
  \midrule
  % ── Medical ──
  AUROC          & Area Under ROC Curve           & Med.    & MedAgentBoard     & 0--1       & 4 \\
  AUPRC          & Area Under PR Curve            & Med.    & MedAgentBoard     & 0--1       & 4 \\
  ROUGE-L        & ROUGE-L (LCS)                  & Med.    & MedAgentBoard     & 0--1       & 4 \\
  BLEU           & BLEU Score                     & Med.    & MedAide           & 0--1       & 3 \\
  SARI           & SARI Score                     & Med.    & MedAgentBoard     & 0--1       & 3 \\
  \midrule
  % ── Framework ──
  T              & Trustworthiness Index          & FW      & TRiSM             & 0--1       & 3 \\
  M\_comp        & Composite Eval.\ Metric        & FW      & TRiSM             & 0--1       & 2 \\
  NDCG@8         & NDCG at 8                      & FW      & ARCANE            & 0--1       & 3 \\
  Goal-Drift     & Goal-Drift Score               & FW      & Shukla 2025       & 0--1       & 2 \\
  Harm-Reduction & Harm Reduction Index           & FW      & Shukla 2025       & 0--1       & 2 \\
  BAF            & Bias Amplification Factor      & FW      & EmergBias         & $>0$       & 3 \\
  ASE            & Agent Scaling Efficiency       & FW      & Survey            & $>0$       & 3 \\
  \midrule
  % ── Operational ──
  TTFT           & Time to First Token            & Ops.    & NVIDIA/MLPerf     & ms$\downarrow$ & 4 \\
  TPOT           & Time Per Output Token          & Ops.    & NVIDIA/MLPerf     & ms$\downarrow$ & 3 \\
  CPT            & Cost Per Task                  & Ops.    & Industry          & USD$\downarrow$ & 4 \\
  TBU            & Token Budget Utilisation       & Ops.    & Survey            & 0--100\%   & 3 \\
  Goodput        & Goodput                        & Ops.    & Databricks        & 0--100\%   & 3 \\
  RunVar         & Run-to-Run Variance            & Ops.    & MAESTRO           & $\sigma^2\downarrow$ & 3 \\
  \midrule
  % ── Tool Use ──
  TSA            & Tool Selection Accuracy        & Tool    & Survey            & 0--100\%   & 4 \\
  TPA            & Tool Parameter Accuracy        & Tool    & Survey            & 0--100\%   & 3 \\
  TCR            & Tool Call Success Rate         & Tool    & Survey            & 0--100\%   & 4 \\
\end{longtable}

\end{document}
